\section{Ethical considerations}

\stitle{Data and privacy.} \name generates its evaluation data entirely from scratch: latent relational sources, tables, and questions are synthesized by LLMs from a domain specification and contain only fictional entities and randomly sampled numerical values. The generated instances therefore include no personally identifying information and no real individuals, and we collect no scraped or user-contributed content. The real tables used in our validation experiments (\S\ref{app:generating_non_rel}) come from public databases used exclusively for evaluation (Northwind, BikeStores, SQL Bootcamp) and from existing research benchmarks (Spider, BIRD, BEAVER, TAT-QA, GRI-QA) built from public corporate and web sources. Where such sources mention named public figures, this information is already public and is neither augmented nor inferred by \name.

\stitle{Use of existing artifacts.} We use all datasets and models in a manner consistent with their intended research use, and we report the license of each artifact in \S\ref{app:license}. Our use is non-commercial and research-only, and any derivatives we produce remain within a research context. The crawled sample databases are used solely for testing, without redistribution (\S\ref{app:generating_non_rel}).

\stitle{Risks and broader impact.} \name is a diagnostic tool for measuring the robustness of Table QA systems, and we see limited potential for direct misuse. The principal risk is interpretive: because the benchmark is synthetic and controlled, strong performance on \name should not be read as a guarantee of real-world reliability, nor should low scores be taken as evidence that a model is unusable in practice. We therefore position \name as a complement to, rather than a replacement for, manually curated benchmarks built from naturally occurring documents (\S\ref{sec:limitations}). Used appropriately, surfacing the multi-table, non-relational, and unit-heterogeneity failure modes, we believe \name can support the development of more reliable analytical interfaces over tabular data, including in financial and environmental reporting settings where silent numerical errors carry real consequences.

\stitle{Human verification and AI assistants.} The correctness study in \autoref{tab:correctness} was carried out by the authors and involves no external annotators or human subjects, raising no participant-compensation or consent concerns. Our use of AI assistants for writing is described in \S\ref{sec:ai}, and the role of LLMs within the generation pipeline is described in \S\ref{sec:approach}.